\scanpage{8}
choisi --- si on \deleted{multiplie} remplace \unsure{$g$|$u$} par $\operatorname{int}(\lambda)u$ avec $\lambda=l_0^\alpha\in L_0$ ($\alpha\in\hat{\mathbb Z}^{\unsure{*}}$), alors $g$ est remplacé par
\[
\lambda g\rho(\unsure{l_0|\lambda})^{-1}=\uline{l_0^\alpha gl_1^{-\alpha}}.
\]
Ceci posé, \deleted{\unsure{les}} les relations $(*)$ et $(**)$ déterminent $u$ en termes de $p\in\hat{\mathbb Z}^{*}$ et de $g\in\pi$, par
\[
\left\{\begin{aligned}
u(l_0)&=l_0^p\\
u(l_1)&=u\rho(l_0)=\operatorname{int}(g)\rho\underbrace{u(l_0)}_{l_0^p}=\operatorname{int}(g)l_1^p\\
u(l_\infty)&=u\rho(l_1)=\operatorname{int}(g)\rho\underbrace{u(l_1)}_{\operatorname{int}(g)l_1^p}\\
&=\operatorname{int}(g)\underbrace{\rho\operatorname{int}(g)\rho^{-1}}_{\operatorname{int}(\rho(g))}\underbrace{(\rho l_1^p)}_{l_\infty^p}
=\deleted{\unsure{\operatorname{int}g\rho(g)l_\infty^p}}\inserted{\operatorname{int}(g\rho(g))l_\infty^p}
\end{aligned}\right.
\]
\margin{\uline{NB} On voit que $u$ automatiquement respecte la structure à lacets, avec le multiplicateur $p$}

Il faut vérifier, pour qu'on ait bien un endom. de $\pi$, la relation $l_0'l_1'l_\infty'=1$ i.e.
\[
\boxed{l_0^p\cdot\operatorname{int}(g)l_1^p\cdot\operatorname{int}(g\rho(g))l_\infty^p=1}
\]
ou encore
\[
l_0^p(gl_1^pg^{-1})(g\rho(g)l_\infty^p\rho(g)^{-1}g^{-1})=1
\]
i.e.
\[
l_0^pgl_1^p\rho(g)l_\infty^p\rho(g)^{-1}g^{-1}=1
\]
soit encore, en faisant passer le $g^{-1}$ à gauche
\[
\deleted{l_0^p\cdot\operatorname{int}(g)\bigl(l_1^p\cdot\operatorname{int}(\rho(g))(l_\infty^p)\bigr)=1}
\]
\[
\operatorname{int}(g^{-1})l_0^p\cdot l_1^p\cdot\operatorname{int}(\rho(g))(l_\infty^p)=1
\]
Les relations $(*)\unsure{,}$ et $(**)$ sont vérifiées par construction (pour $**$, il suffit $\unsure{=|\bar m}$ de la vérifier pour la valeur \unsure{aux|des} deux mb en $l_0,l_1$).
\mspage{7}
